% Source: https://artofproblemsolving.com/community/c3h2168634_perimeter_of_triangle
% Stable identity: aops-topic-2168634-q01
The two tangent segments from an external point to a circle are equal. From $S$ this gives $SU=SQ$, and from $T$ it gives $TU=TR$.

\begin{asy}
usepackage("amsmath");
size(6cm);
defaultpen(fontsize(10pt));
pair O=(0,0), P=(-9,0);
real t=aCos(5/9);
pair Q=5*dir(180-t), R=5*dir(180+t), U=5*dir(166);
pair d=dir(76);
pair Sp=extension(U,U+d,P,Q), Tp=extension(U,U+d,P,R);
draw(circle(O,5));
draw(P--Q,linewidth(0.8)); draw(P--R,linewidth(0.8));
draw(Sp--Tp,blue+linewidth(1));
draw(O--Q,dashed); draw(O--P,dashed); draw(O--U,dashed);
draw(P--Sp--Tp--cycle,red+linewidth(1.2));
dot(O); dot(P); dot(Q); dot(R); dot(U); dot(Sp); dot(Tp);
label("$O$",O,E); label("$P$",P,W); label("$Q$",Q,N); label("$R$",R,S);
label("$U$",U,NE); label("$S$",Sp,NW); label("$T$",Tp,SW);
label("$9$",(-7,0),S); label("$5$",(O+Q)/2,NE);
label("\mbox{}",(min(currentpicture,user=true)+max(currentpicture,user=true))/2,black); // keeps the text after the figure black when LaTeX typesets the labels inline
\end{asy}

Since $U$ lies between $S$ and $T$,
\[PS+ST+TP=PS+SU+UT+TP=(PS+SQ)+(PT+TR)=PQ+PR .\]
The radius $OQ$ is perpendicular to the tangent $PQ$, so by the Pythagorean theorem in right triangle $OQP$,
\[PQ=\sqrt{PO^2-OQ^2}=\sqrt{81-25}=\sqrt{56}=2\sqrt{14},\]
and $PR=PQ$. The perimeter of triangle $PST$ is $2PQ=\boxed{4\sqrt{14}}$, whatever the position of $U$ on the arc.
